\documentclass[10pt]{article}
\usepackage{lingmacros}
\usepackage{tree-dvips}
\usepackage{makecell}
\usepackage{xfrac}
\usepackage{multicol}
\usepackage{graphicx}
\usepackage{siunitx}
\usepackage{hyperref}
\usepackage{wrapfig}
\usepackage[T1]{fontenc}
\usepackage[a4paper, margin=3cm]{geometry}
\usepackage{multirow}
\usepackage{titlesec}
\begin{document}
\tableofcontents
\pagebreak
\section{}
\textit{Conservazione di salute e benessere di pazienti e operatori, riducendo i rischi derivati dall'uso di \textbf{radiazioni ionizzanti} in attività giustificate da \textbf{benefici} garantiti.}.\\
\textit{Le radiazioni sono emissione e propagazione di \textbf{energia} attraverso spazio e materia a livello microscopico}, queste sono dovute al passaggio dallo stato eccitato a quello basale degli elettroni che circondano i nuclei. Le radiazioni non ionizzanti non sono in grado di interagire con i tessuti in modo lesivo ($<$12eV), mentre quelle \underline{ionizzanti} sì e si classificano in direttamente (particelle $\alpha$ e $\beta$, corpuscolate) e indirettamente (raggi X e $\gamma$, elettromagnetiche) ionizzanti. \\
Le radiazioni possono inoltre avere origine naturale in nuclei con n. atomico 81-92 (69\% come radon, raggi cosmici, terreno e interne) o artificiale da bombardamento (medicina nucleare, prodotti di consumo, RX), queste possono essere prodotte da macchine radiogene attive (flusso intenso e irradiazione esterna) o da sostanze radioattive costantemente (da sorgente sigillata se formate di materiale solido inattivo o inattivate da apposito involucro o non sigillate).\\
La \textbf{radioattività} è definita come l'\textit{insieme di fenomeni fisico-nucleari (decadimento) con cui un nuclide (atomo con numero di protoni e elettroni definito) instabile cede energia come radiazioni per raggiungere lo stato di stabilità, in un lasso di tempo chiamato \textbf{tempo di decadimento}}.\\
La \underline{vita media} di un isotopo è il tempo che impiega per trasformarsi in altri, mentre il \underline{tempo di dimezzamento} è il tempo che impiegano la metà degli atomi radioattivi a trasformarsi spontaneamente. \\
Per limitare l'interazione delle radiazioni con la materia è possibile usare schermature e DPI, aumentare la distanza (dose inversamente proporzionale al quadrato della distanza) e ridurre il tempo; le radiazioni possono provenire da radionuclidi manipolati, pazienti iniettati, sorgenti di calibrazione e rifiuti contaminati, il tipo di contaminazione esterna più frequente è tramite aria, acqua, cibo, superfici di lavoro e radiofarmaci.\\
La \textbf{contaminazione interna} avviene per contatto diretto tramite le mani e la pelle, una contaminazione ridotta ma ripetuta nel tempo risulta dannosa a lungo andare ed è per questo necessario contenere quanto più possibile le sorgenti non sigillate durante la manipolazione: le \underline{sorgenti liquide} devono essere contenute in doppio contenitore integro circondati da carta assorbente e da un foglio di plastica e le \underline{sorgenti gassose} (I$^{131}$) devono essere usate sotto cappa aspirante schermata. \\
L'\textbf{esposizione esterna} è dovuta a irradiazione da una sorgente di radiazioni X, $\gamma$ o $\beta$ in prossimità del soggetto, le prime due rappresentano un rischio molto elevato per il loro alto potere di penetrazione in particolare per gonadi e midollo. Le radiazioni $\alpha$ sono pericoloso solamente per l'occhio nel caso di esposizioni ravvicinate, le $\beta$ sono legate invece ad un contatto diretto o a breve distanza. I raggi X possono essere schermati da coperture in metallo a differenza dei raggi $\gamma$ che le attraversano. Tutte interagiscono con la materia perché assorbite 
\end{document}